\documentclass[10pt, letterpaper]{article}

% Packages:
\usepackage[
    ignoreheadfoot, % set margins without considering header and footer
    top=1.4 cm, % seperation between body and page edge from the top
    bottom=1.4 cm, % seperation between body and page edge from the bottom
    left=1.5 cm, % seperation between body and page edge from the left
    right=1.5 cm, % seperation between body and page edge from the right
    footskip=1.0 cm, % seperation between body and footer
    % showframe % for debugging 
]{geometry} % for adjusting page geometry
\usepackage{titlesec} % for customizing section titles
\usepackage{tabularx} % for making tables with fixed width columns
\usepackage{array} % tabularx requires this
\usepackage[dvipsnames]{xcolor} % for coloring text
\definecolor{primaryColor}{RGB}{0, 0, 0} % define primary color
\usepackage{enumitem} % for customizing lists
\usepackage{fontawesome5} % for using icons
\usepackage{amsmath} % for math
\usepackage[
    pdfauthor={Curtis Welter},
    pdfcreator={LaTeX},
    colorlinks=true,
    urlcolor=primaryColor
]{hyperref} % for links, metadata and bookmarks
\usepackage[pscoord]{eso-pic} % for floating text on the page
\usepackage{calc} % for calculating lengths
\usepackage{bookmark} % for bookmarks
\usepackage{lastpage} % for getting the total number of pages
\usepackage{changepage} % for one column entries (adjustwidth environment)
\usepackage{paracol} % for two and three column entries
\usepackage{ifthen} % for conditional statements
\usepackage{needspace} % for avoiding page brake right after the section title
\usepackage{iftex} % check if engine is pdflatex, xetex or luatex

% Ensure that generate pdf is machine readable/ATS parsable:
\ifPDFTeX
    \input{glyphtounicode}
    \pdfgentounicode=1
    \usepackage[T1]{fontenc}
    \usepackage[utf8]{inputenc}
    \usepackage{lmodern}
\fi

\usepackage{charter}

% Some settings:
\raggedright
\AtBeginEnvironment{adjustwidth}{\partopsep0pt} % remove space before adjustwidth environment
\pagestyle{empty} % no header or footer
\setcounter{secnumdepth}{0} % no section numbering
\setlength{\parindent}{0pt} % no indentation
\setlength{\topskip}{0pt} % no top skip
\setlength{\columnsep}{0.15cm} % set column seperation
\pagenumbering{gobble} % no page numbering

\titleformat{\section}{\needspace{2\baselineskip}\bfseries\large}{}{0pt}{}[\vspace{1pt}\titlerule]

\titlespacing{\section}{
    % left space:
    -1pt
}{
    % top space:
    0.3 cm
}{
    % bottom space:
    0.2 cm
} % section title spacing

\renewcommand\labelitemi{$\vcenter{\hbox{\small$\bullet$}}$} % custom bullet points
\newenvironment{highlights}{
    \begin{itemize}[
        topsep=0.10 cm,
        parsep=0.10 cm,
        partopsep=0pt,
        itemsep=0pt,
        leftmargin=0 cm + 10pt
    ]
}{
    \end{itemize}
} % new environment for highlights


\newenvironment{highlightsforbulletentries}{
    \begin{itemize}[
        topsep=0.10 cm,
        parsep=0.10 cm,
        partopsep=0pt,
        itemsep=0pt,
        leftmargin=10pt
    ]
}{
    \end{itemize}
} % new environment for highlights for bullet entries

\newenvironment{onecolentry}{
    \begin{adjustwidth}{
        0 cm + 0.00001 cm
    }{
        0 cm + 0.00001 cm
    }
}{
    \end{adjustwidth}
} % new environment for one column entries

\newenvironment{twocolentry}[2][]{
    \onecolentry
    \def\secondColumn{#2}
    \setcolumnwidth{\fill, 4.5 cm}
    \begin{paracol}{2}
}{
    \switchcolumn \raggedleft \secondColumn
    \end{paracol}
    \endonecolentry
} % new environment for two column entries

\newenvironment{threecolentry}[3][]{
    \onecolentry
    \def\thirdColumn{#3}
    \setcolumnwidth{, \fill, 4.5 cm}
    \begin{paracol}{3}
    {\raggedright #2} \switchcolumn
}{
    \switchcolumn \raggedleft \thirdColumn
    \end{paracol}
    \endonecolentry
} % new environment for three column entries

\newenvironment{header}{
    \setlength{\topsep}{0pt}\par\kern\topsep\centering\linespread{1.5}
}{
    \par\kern\topsep
} % new environment for the header

\newcommand{\placelastupdatedtext}{% \placetextbox{<horizontal pos>}{<vertical pos>}{<stuff>}
  \AddToShipoutPictureFG*{% Add <stuff> to current page foreground
    \put(
        \LenToUnit{\paperwidth-2 cm-0 cm+0.05cm},
        \LenToUnit{\paperheight-1.0 cm}
    ){\vtop{{\null}\makebox[0pt][c]{
        \small\color{gray}\textit{Last updated in September 2024}\hspace{\widthof{Last updated in September 2024}}
    }}}%
  }%
}%

% save the original href command in a new command:
\let\hrefWithoutArrow\href

% new command for external links:


\begin{document}
    \newcommand{\AND}{\unskip
        \cleaders\copy\ANDbox\hskip\wd\ANDbox
        \ignorespaces
    }
    \newsavebox\ANDbox
    \sbox\ANDbox{$|$}

    \begin{header}
        \fontsize{25 pt}{25 pt}\selectfont Curtis Welter

        \vspace{5 pt}

        \normalsize
        \mbox{Atlanta, GA 30316}%
        \kern 5.0 pt%
        \AND%
        \kern 5.0 pt%
        \mbox{\hrefWithoutArrow{mailto:weltercurtis@gmail.com}{weltercurtis@gmail.com}}%
        \kern 5.0 pt%
        \AND%
        \kern 5.0 pt%
        \mbox{\hrefWithoutArrow{tel:+1-678-382-4977}{678 382 4977}}%
        \kern 5.0 pt%
        \AND%
        \kern 5.0 pt%
        \mbox{\hrefWithoutArrow{https://linkedin.com/in/curtis-welter}{linkedin.com/in/curtis-welter}}%
        \kern 5.0 pt%
    \end{header}

    \vspace{5 pt - 0.3 cm}

    \section{Education}
        
        \begin{twocolentry}{
            01/2026 – Present
        }
            \textbf{Georgia Institute of Technology}, MS in Computer Science-Expected Graduation 2028\end{twocolentry}

        \vspace{0.10 cm}
        \begin{onecolentry}
            \begin{highlights}
                \item \textbf{Coursework:} Operating Systems, Cognitive Science, AI Techniques in Robotics, Computer Graphics
                \item Specialization in Computational Perception and Robotics
            \end{highlights}
        \end{onecolentry}
        \vspace{0.20 cm}
        \begin{twocolentry}{
            08/2018 – 05/2022
        }
            \textbf{Georgia Institute of Technology}, BS in Computer Science\end{twocolentry}

        \vspace{0.10 cm}
        \begin{onecolentry}
            \begin{highlights}
                \item \textbf{Coursework:} Computer Systems \& Networks, Information Security, Natural Language Processing, Computer Vision, Machine Learning, Artificial Intelligence and Robotics
                \item Studied Abroad Summer 2019 at Barcelona School of Informatics
            \end{highlights}
        \end{onecolentry}

    \section{Languages and Technologies}

        \begin{onecolentry}
            \begin{highlightsforbulletentries}
            
            \item Python, Java, C++, C\#, Kotlin, Swift, Javascript, C, R, MATLAB, Microsoft SQL Server, Assembly
    
            \item React Native, Angular; Azure, .NET, AWS, GCP; UNIX/Linux, Android, iOS; Git, Docker; PyTorch; LaTeX; AI/ML
    
            \item Adobe Photoshop, Premiere, After Effects, Illustrator; Blender; Unreal Engine, Unity, Godot; Fusion 360;
    
            \end{highlightsforbulletentries}
        \end{onecolentry}

    
    \section{Experience}
        \vspace{0.05 cm}
        \begin{twocolentry}{
            11/2023 - 03/2026
        }
            \textbf{Centers for Disease Control}, Full Stack Developer\end{twocolentry}

        \vspace{0.10 cm}
        \begin{onecolentry}
            \begin{highlights}
                \item Second year developer on the New Data Source Development Team in support of the National Syndromic Surveillance Program. Creating new pipelines utilizing AI/ML techniques for the intake, analysis, and visualization of public healthcare big data from a variety of sources to monitor and predict national health emergencies. \textit{Python, Java, MS SQL}
                \item First year worked in support of the National Center for Chronic Disease Prevention and Health Promotion's suite of mobile apps, added new functionality and security features to protect and validate user data. Developed on a variety of frameworks and platforms including android, iOS, and web. \textit{Kotlin, .NET, C\#, MySQL, Swift, Ionic, Javascript}
            \end{highlights}
        \end{onecolentry}


        \vspace{0.2 cm}

        \begin{twocolentry}{
            03/2023 - 09/2023
        }
            \textbf{Novarata}, Full Stack Developer\end{twocolentry}

        \vspace{0.1 cm}
        \begin{onecolentry}
            \begin{highlights}
                \item Assisted in development of kernel level android driver for smart watch’s network radio, increasing battery life of device from 20 to 50 hours. \textit{Java and C}
                \item Added new features to client company’s in house industrial ERP and CRM systems. Added features to frontend to allow for more customizable quotes, and resolved email API bugs. \textit{Angular, C\#, .NET, and MS SQL}
            \end{highlights}
        \end{onecolentry}

        \vspace{0.2 cm}

        \begin{twocolentry}{
            01/2020 - 05/2022
        }
            \textbf{Georgia Tech Research Institute}, Research Assistant\end{twocolentry}

        \vspace{0.1 cm}
        
        \begin{onecolentry}
            \begin{highlights}
                \item Conducted guided research on security vulnerabilities within Intel C++ Compiler’s Parallel Processing Optimization. Authored section of official threat assessment report detailing findings in accordance with NIST standard. \textit{C++}
                \item Developed document analysis utility on contract with the Department of Defense utilizing Computer Vision and Natural Language Processing techniques. \textit{Python}
                \item Assisted in development of a naval ship tracking and networking system. Created python utility to parse, normalize, and visualize location data reported by a variety of seafaring vessels. \textit{Python}
            \end{highlights}
        \end{onecolentry}

        
    \section{Projects}
        \begin{highlights}
            \item \textbf{Read-A-Rhyme} (2022) Capstone project developing a mobile app to teach literacy to elementary school children using nursery rhymes. Project team of 5 in coordination with early literacy researchers at our client Hearatale. \textit{React Native}
            \item \textbf{Peach Picking Robot} (2021) Joined a team in Georgia Tech’s VIP program to develop an agricultural robot. Implemented computer vision for ripe peach detection and a denoising algorithm for LiDAR navigation. \textit{Python, C++}
            \item \textbf{Moog Hackathon} (2020) Soldered and programmed arduino based microcontroller to augment the capabilities of a Moog ‘Werkstatt’ Analog Synthesizer musical instrument, adding MIDI keyboard support. \textit{C (language)}
            \item \textbf{HackGT Hackathon} (2018-2019) 2 year participant. 1st year my team built an app which identifies comic book characters using machine learning and cloud based computing. 2nd Year won award for Best Use of Google Cloud API for litter tracking and cleanup app with Google Maps and Vision integration. \textit{Android, Java, Kotlin}
            \item \textbf{Ludum Dare Game Jam} (2014-Present) Regular participant in game jams as part of a team and as a solo developer. Current platform NES/Gameboy Homebrew. \textit{Actionscript 3/Flash, C\#, Unity, C++, Unreal Engine, Assembly}
        \end{highlights}
\end{document}